% ---------------------------------------------------------------
%  Muc II.6 -- Huong phat trien
% ---------------------------------------------------------------
\subsection{Hướng phát triển tiếp theo}

Phần đã hoàn thành và đang được sử dụng là vòng chạy phục vụ \textbf{học sinh}:
tạo đề, làm bài, chấm điểm, trả kết quả về Google Classroom, cùng trang thống
kê cơ bản cho giáo viên. Các hướng phát triển dưới đây tập trung làm sâu
\textbf{phần dành cho giáo viên} --- phần mà theo tôi mới là chỗ sáng kiến này
phát huy giá trị lớn nhất đối với tổ chuyên môn.

\subsubsection{Khu làm việc riêng cho giáo viên}

Đây là hướng ưu tiên số một, gồm bốn chức năng:

\begin{enumerate}[label=(\arabic*), leftmargin=1.2cm, itemsep=3pt]
  \item \textbf{Tự đặt ma trận chi tiết cho từng đề.} Hiện ma trận lấy theo
        bảng cấu hình chung của loại bài kiểm tra. Giáo viên cần đặt được ma
        trận riêng cho từng đề: bài nào mấy câu, mức độ nào mấy câu, dạng câu
        nào --- tức là đưa chính bảng ma trận quen thuộc của tổ chuyên môn lên
        màn hình.
  \item \textbf{Sinh một lượt nhiều mã đề.} Thay vì bấm tạo nhiều lần, giáo
        viên nhập số mã đề cần có và nhận về một tệp gồm đủ các mã đề cùng cấu
        trúc, khác số liệu, kèm bản lời giải và bảng đáp án tổng hợp của tất cả
        các mã --- dùng in kiểm tra trực tiếp.
  \item \textbf{Tải mã nguồn \LaTeX\ của đề về máy.} Giáo viên lấy được tệp
        \texttt{.tex} của đề để tự chỉnh sửa theo ý mình: thêm một câu riêng,
        đổi cách trình bày, ghép vào mẫu đề của trường. Đây là yêu cầu thực tế
        của giáo viên quen dùng \LaTeX\ --- hệ thống lo phần nặng nhọc là chọn
        câu và sinh số liệu, còn quyền biên tập cuối cùng vẫn thuộc về giáo
        viên.
  \item \textbf{Quản lí lớp và kho đề đã giao.} Danh sách lớp, danh sách học
        sinh, các đề đã giao cho từng lớp, kết quả từng đợt --- để giáo viên
        không phải nhớ mình đã cho lớp nào làm đề nào.
\end{enumerate}

\ghichu{Bốn chức năng này em sẽ làm ngay sau khi bài sáng kiến hoàn thành, theo
đúng yêu cầu của chị: chỉ ra đề và lấy mã nguồn, không cần làm bài trực tuyến ở
nhánh này.}

\subsubsection{Mở rộng ngân hàng câu hỏi}

Ngân hàng hiện mới phủ chương I lớp 10. Kế hoạch mở rộng theo thứ tự: hoàn
thiện các chương còn lại của lớp 10, sau đó đến lớp 11 và lớp 12. Riêng dạng
câu trả lời ngắn mức vận dụng cao đang là ô mỏng nhất của ma trận --- đây là
việc cần làm trước tiên, và dòng ``[THIẾU CÂU HỎI]'' trong đề chính là danh
sách việc cần làm do hệ thống tự chỉ ra.

Về lâu dài, nếu được tổ chuyên môn hưởng ứng, mỗi giáo viên đóng góp hàm sinh
câu hỏi cho một vài bài mình dạy kĩ nhất thì ngân hàng chung của tổ sẽ lớn rất
nhanh --- đây là mô hình làm việc mà tôi mong muốn hướng tới.

\subsubsection{Tổ chức lại ngân hàng khi quy mô lớn lên}

Khi ngân hàng vượt vài nghìn câu hỏi, việc tra cứu và chống trùng lặp sẽ thành
vấn đề mới. Hướng xử lí dự kiến chia làm hai mức:

\begin{itemize}[leftmargin=1.2cm, itemsep=3pt]
  \item \textbf{Mức 1 --- chỉ mục ngân hàng.} Một chương trình quét toàn bộ
        thư mục câu hỏi và lập ra bảng tổng hợp: bài nào, mức độ nào đang
        thiếu câu; mã nào có trong bảng ánh xạ nhưng chưa có hàm sinh; mỗi bài
        hiện có bao nhiêu câu mỗi mức. Kèm theo là bộ dò các câu hỏi trùng
        nghĩa. Mức này không cần công nghệ gì mới và là việc làm trước.
  \item \textbf{Mức 2 --- tìm kiếm theo nghĩa.} Khi giáo viên cần tìm câu hỏi
        mà không nhớ mã, có thể dùng kĩ thuật \textit{nhúng vecto}: mỗi câu hỏi
        được biểu diễn thành một vecto trong không gian nhiều chiều sao cho hai
        câu gần nhau về nghĩa thì hai vecto gần nhau, và việc tìm kiếm trở
        thành tìm vecto có cô-sin góc lớn nhất. Cơ sở dữ liệu đang dùng đã hỗ
        trợ sẵn kiểu dữ liệu này nên không phải dựng thêm hạ tầng.
\end{itemize}

Cần nói rõ để tránh hiểu nhầm: kĩ thuật vecto \textbf{không} thay thế cách chọn
câu hiện tại. Việc chọn câu theo ma trận đòi hỏi đúng tuyệt đối về bài, mức độ
và loại câu, nên phải tra theo mã định danh --- đó là tra từ điển, chính xác
hoàn toàn. Tìm kiếm vecto chỉ cho kết quả gần đúng có xếp hạng, phù hợp với
việc \textit{tìm}, không phù hợp với việc \textit{ra đề}. Vì vậy nếu triển
khai, nó sẽ là một lớp tìm kiếm phụ đặt bên cạnh cho giáo viên dùng, chứ không
nằm trong luồng sinh đề.

\subsubsection{Chấm câu tự luận}

Hiện phần tự luận 3 điểm học sinh làm ra giấy. Hướng phát triển là cho học sinh
chụp ảnh bài làm viết tay và bài chấm được đối chiếu với lời giải chuẩn mà hàm
sinh câu hỏi đã tạo ra. Phần kĩ thuật đã dựng xong và chạy thông với ảnh thử,
nhưng tôi chủ động chưa đưa vào sử dụng vì chưa kiểm chứng được độ chính xác
khi đọc chữ viết tay thật của học sinh --- một chức năng chấm điểm mà chưa chắc
chắn thì chưa nên dùng. Màn hình hiện đang ghi rõ ``chức năng chấm tự luận đang
được xây dựng'' để học sinh không hiểu nhầm.

\subsubsection{Phân tích kết quả và gợi ý luyện tập cá nhân hoá}

Hệ thống đã lưu đủ dữ liệu để biết mỗi học sinh sai nhiều ở bài nào, mức độ
nào. Bước tiếp theo là tự động chỉ ra điểm yếu của từng em và gợi ý ngay một đề
luyện tập đúng vào phần yếu đó, chỉ bằng một nút bấm. Khi có chức năng này,
vòng ``học --- kiểm tra --- biết mình yếu chỗ nào --- luyện đúng chỗ đó'' mới
thực sự khép kín, và đúng tinh thần \textit{đánh giá vì sự tiến bộ của người
học} của Thông tư 22/2021/TT-BGDĐT.

\subsubsection{Hoàn thiện phần tài khoản giáo viên}

Hiện việc đăng kí tài khoản mới trên hệ thống dành cho học sinh; giáo viên dùng
các chức năng riêng qua đường dẫn trực tiếp. Cần bổ sung luồng đăng kí và phân
quyền cho giáo viên để nhiều giáo viên trong tổ cùng dùng độc lập, mỗi người
quản lí lớp của mình.

\clearpage
